\lecture{15}{Tue 26 Nov 2024 09:30}{Ideals and Homomorphisms}
\begin{definition}[Maximal Ideal]\label{dfn:MID}
	An ideal $M$ in a ring $R$ is maximal if $M\neq R$ and the only ideal strictly containing $M$ is $R$.
\end{definition}

Rings can have multiple maximal ideals. For example,  $p\mathbb{Z}$ is a maximal ideal of $\mathbb{Z}$ for all prime $p$. Moreovr, Krull's theorem states maximal ideals always exist.

\begin{theorem}\label{thm:akj}
	Let $R$ be a commutative ring with unity, and let $M\subset R$ be an ideal. Then $M$ is maximal if and only if $R /M$ is a field.
\end{theorem}
\begin{proof}
	done in homework 10
\end{proof}

Note that $\mathbb{Z} / p\mathbb{Z}$ is a field, proving our above statement that $\mathbb{Z}$ has multiple maximal ideals.

\begin{definition}[Prime Ideal]\label{dfn:PrI}
	A proper ideal $P$ of a commutative ring $R$ is \textbf{prime} if $ab\in P$ if and only if $a\in P$ or $b\in P$.
\end{definition}

An important property of prime numbers is that for a prime $p$, if $p|ab$, then $p|a$ or $p|b$. This is the same idea, but now with ideals instead of like numbers.

\begin{proposition}
	Let $R$ be a commutative ring with unity. Then $P$ is a prime ideal of $R$ if and only if $R /P$ is an integral domain.
\end{proposition}
\begin{proof}
	also done in homework 10
\end{proof}
\begin{corollary}
	Any maximal ideal is prime, but prime ideals are not necessarily maximal.
\end{corollary}
\begin{proof}
	Any field is an integral domain, so it follows immediately from the previous theorem and proposition.
\end{proof}

\begin{corollary}
	A commutative ring with unity ring $R$ is a field if and only if $\left\langle 0 \right\rangle $ is a maximal ideal.
\end{corollary}
\begin{proof}
	done in the homework, except this is not a corollary in a homework but a lemma to the previous theorems.
\end{proof}
\begin{proposition}
	A commutative ring with unity $R$ is an integral domain if and only if $\left\langle 0 \right\rangle $ is a prime ideal.
\end{proposition}
\begin{proof}
	Let $ab=0$. Then either $a$ or $b$ is $0$. Therefore, $R$ is an integral domain.
\end{proof}
\chapter{Ring Homomorphisms}
\begin{definition}[Ring Homomorphism]\label{dfn:RingHomo}
	Let $R,S$ be rings, and let $\phi : R \to S$. Then $\phi $ is a ring homomorphism if and only if for all $a,b\in S$,
	\[
		\phi (ab)=\phi (a)\phi (b)\qquad \text{and}\qquad \phi (a+b)=\phi (a)+\phi (b)
	.\]
\end{definition}
\begin{definition}[Ring Isomorphism]\label{dfn:RingIso}
	Let $R,S$ be rings and let $\phi : R \to S$ be a ring homomorphism. Then $\phi $ is an isomorphism if and only if it is bijective. We write $R\cong S$ to denote that $R$ is isomorphic to $S$.
\end{definition}
\begin{definition}[Kernel]\label{dfn:RingKer}
	Let $R,S$ be rings and let $\phi : R \to S$ be a ring homomorphism. Define the kernel of $\phi $ as the set
	\[
		\ker \phi =\left\{ r\in R \mid \phi (r)=0 \right\} 
	.\]
\end{definition}
\begin{proposition}
	Let $R,S$ be rings and let $\phi : R \to S$ be a ring homomorphism. Then $\ker \phi $ is an ideal of $R$.
\end{proposition}
\begin{proof}
	Done in homework. I will take that $\ker \phi $ is a ring to be trivial. Let $r\in \ker \phi $. Note that
	\[
		\phi (ar)=\phi (a)0=0
	.\]
	Therefore, $ar\in \ker \phi $.
\end{proof}
\begin{proposition}
	Let $R$ be a ring with $I$ an ideal. Define $\phi : R \to R /I$ as the map $r\mapsto r+I$. Then $\ker \phi =I$.
\end{proposition}
\begin{proof}
	Let $r\in \ker \phi $. It follows that $\phi (r)=r+I=I$. By properties of cosets, $r\in I$. Now let $r\in I$. It follows by properties of cosets that $\phi (r)=r+I=I$. Therefore, $r\in \ker \phi $.
\end{proof}
\begin{remark}
	We call $\phi : R \to R/I$ defined as $r\mapsto r+I$ the \textbf{canonical} \textbf{homomorphism}.
\end{remark}
Here's an interesting ring homomorphism. Let $\phi_\alpha $ be a linear functional. That is, let $\phi_\alpha : C[a,b] \to \mathbb{R}$ be the map $f\mapsto f(\alpha )$. Then  $\phi_\alpha $ is a ring homomorphism for all $\alpha \in\mathbb{R}$.

\begin{theorem}[Properties]\label{thm:ka}
	Let $\phi : R \to S$ be a ring homomorphism.
	\begin{itemize}
		\item If $R$ is a commutative ring, $\phi (R)$ is a commutative ring.
		\item $\phi (0)=0$.
		\item Let $R$ be a ring with unity. If $\phi $ is surjective, then $\phi (1)=1$.
		\item If $R$ is a field and $\phi (R)\neq \left\{ 0 \right\}  $, then $\phi (R)$ is a field.
	\end{itemize}
\end{theorem}
\begin{proof}
	trivial.
\end{proof}
\begin{theorem}[First Isomorphism Theorem]\label{thm:FIT}
	Let $R,S$ be rings and let $\psi  : R \to S$ be a ring homomorphism. Then the map  $\eta  : R /\ker \psi   \to \psi  (R) $ defined as $r+\ker \psi  \mapsto \psi  (r)$ is a unique isomorphism. Moreover, if $\phi : R \to R /\ker \psi $ is the canonical isomorphism theorem, then $\psi =\eta \circ \phi $.
\end{theorem}
\begin{theorem}[Second Isomorphism Theorem]\label{thm:SIT}
	Let $R$ be a ring and let $I,J$ be ideals of $R$. Then
	\[
		I /I\cap J \cong (I+J) /J
	.\]
\end{theorem}
\begin{theorem}[Third Isomorphism Theoremnm]\label{thm:FIT}
	Let $R$ be a ring and let $I,J$ be ideals of $R$ with $J \subset I \subset R$. Then
	\[
		R /I \cong \frac{R /J}{I /J}
	.\]
\end{theorem}
\begin{theorem}[Correspondence Theorem]\label{thm:FOIT}
	Problem 5a of hw 10? idk
\end{theorem}
